% =========================================================================
% main.tex — PhD Thesis LaTeX Template
% -------------------------------------------------------------------------
% Originally developed by Indrajit Ghosh (2024–2025)
% Indian Statistical Institute, Bangalore Centre
%
% License: Creative Commons Attribution-NonCommercial 4.0 International (CC BY-NC 4.0)
% https://creativecommons.org/licenses/by-nc/4.0/
%
% This template provides a clean structure for typesetting PhD theses using LaTeX.
% It includes title setup, PDF metadata, GPG signing guide, and example chapters.
%
% To adapt it for your own use:
%   1. Replace the constants below with your information.
%   2. Edit or add chapters under the "chapters/" directory.
%   3. Compile with: pdflatex → biber → pdflatex → pdflatex
%
% For questions or suggestions, contact: indrajitghosh912@gmail.com
% =========================================================================

% How to Use the \index{} Command
% ------------------------------------------------------------
% The \index{} command allows you to create entries for the index of your document.
% Follow the steps below to enable and use it in your LaTeX project:
%
% 1. Use the \index{} command immediately after the word or phrase you want to index.
%    For example:
%       Hi, I am Indrajit Ghosh, \index{Indrajit} the author of this \LaTeX\ \index{\LaTeX} file.
%
%    In this example:
%       - "Indrajit" will appear as an index entry.
%       - "\LaTeX" will also appear as an index entry.
%
% 2. To generate the index, compile your document as follows:
%       pdflatex main
%       makeindex main
%       pdflatex main
%
% This will produce an index section (usually at the end of your document) listing
% all indexed terms with their corresponding page numbers.
% ------------------------------------------------------------



\documentclass[a4paper,11pt,twoside,openright]{book}

% Input color definitions
% colors.tex
% Author: Indrajit Ghosh
% Created On: Jan 24, 2025
%

\usepackage{xcolor}

\definecolor{royalblue}{HTML}{4169E1}          % Royal Blue
\definecolor{darkgreen}{HTML}{006400}          % Dark Green
\definecolor{crimson}{HTML}{DC143C}            % Crimson
\definecolor{goldenrod}{HTML}{DAA520}          % Goldenrod
\definecolor{darkorange}{RGB}{255,140,0} % RGB for Dark Orange
\definecolor{midnightblue}{RGB}{25, 25, 112}  % Midnight Blue
\definecolor{darkred}{RGB}{139, 0, 0}  % Dark Red
\definecolor{saltgray}{RGB}{198, 198, 198} % A soft gray, resembling salt color
\definecolor{softblue}{HTML}{1A73E8}    % Soft blue for links
\definecolor{deepblue}{HTML}{0056B3}    % Darker blue for URLs
\definecolor{forestgreen}{HTML}{228B22} % Forest green for citations


% These define custom colors for use with \textcolor and other commands
\definecolor{indraRed}{rgb}{0.593, 0.183, 0.183}
\definecolor{indraPink}{rgb}{0.858, 0.188, 0.478}
\definecolor{indraBlue}{rgb}{0, 0.199, 0.398}
\definecolor{madridBlue}{rgb}{0.199, 0.199, 0.695}
\definecolor{metropolisThemeColor}{rgb}{0.105, 0.214, 0.234}
\definecolor{metropolisBarColor}{rgb}{0.984, 0.0.515, 0.015}
\definecolor{UBCblue}{rgb}{0.04706, 0.13725, 0.26667} % UBC Blue (primary)
\definecolor{UBCgrey}{rgb}{0.3686, 0.5255, 0.6235}    % UBC Grey (secondary)


%%%%%%%%%%%%%%%%%%%%%%%%%%%%%%%%%%%%%%%%%%%%%%%%%%%
% Thesis Constants (change as needed)
%%%%%%%%%%%%%%%%%%%%%%%%%%%%%%%%%%%%%%%%%%%%%%%%%%%
\newcommand{\ThesisTitle}{A Comedic Exploration of Algebraic Absurdities: Functorial Adventures in Infinite Dimensions}
\newcommand{\ThesisSubtitle}{(With Occasional Proofs and Existential Crises)} % Keep it empty if not required
\newcommand{\Author}{Indrajit Ghosh}
\newcommand{\ThesisYear}{May 2025}
\newcommand{\AuthorHomepage}{https://indrajitghosh.onrender.com/}
\newcommand{\PhdAdvisor}{Dr. Advisor Here}
\newcommand{\AdvisorHomepage}{https://www.isibang.ac.in/~statmath/faculty.html}
\newcommand{\Institute}{Indian Statistical Institute}
\newcommand{\Department}{Theoretical Statistics \& Mathematics Unit}
\newcommand{\Centre}{Bangalore Centre}

\newcommand{\AuthorAddress}{
    {\small%
        \Department\\%
        \Institute%
        \ifthenelse{\equal{\Centre}{}}{}{, \Centre}%
    }
}

%%%%%%%%%%%%%%%%%%%%%%%%%%%%%%%%%%%%%%%%%%%%%%%%%%%%%%%%%%%%%%%
% PDF Metadata
%%%%%%%%%%%%%%%%%%%%%%%%%%%%%%%%%%%%%%%%%%%%%%%%%%%%%%%%%%%%%%%
\newcommand{\pdfTitle}{Indrajit Ghosh's Mathematical Comedy Thesis 2025}
\newcommand{\pdfAuthor}{Indrajit Ghosh}
\newcommand{\pdfSubject}{Mathematics and Existential Humor}
\newcommand{\pdfKeywords}{Operator Algebras, Humor, Infinite Dimensions, Set Theory}
\newcommand{\pdfCreator}{TeXLive}
\newcommand{\pdfCreationDate}{\today}
\newcommand{\pdfColorLink}{true}
\newcommand{\pdfLinkColor}{royalblue}
\newcommand{\pdfUrlColor}{deepblue}
\newcommand{\pdfCiteColor}{forestgreen}

%%%%%%%%%%%%%%%%%%%%%%%%%%%%%%%%%%%%%%%%%%%%%%%%%%%%%%%%%%%%%%%
% Include preamble and AMS environment setup
%%%%%%%%%%%%%%%%%%%%%%%%%%%%%%%%%%%%%%%%%%%%%%%%%%%%%%%%%%%%%%%
%% preamble.tex
%% This file contains the preambles for the PhD Thesis.
%% Author: Indrajit Ghosh
%% Date: Oct 21, 2024

%%----------------------------------------------------
%%		 Preamble for the project
%%-----------------------------------------------------
 
\usepackage[utf8]{inputenc}
\usepackage{graphicx}
\usepackage{enumitem} % Give extra customization on top of itemize and enumerate
\usepackage{ifthen}

\usepackage{amsmath, amsthm, amssymb} % for various mathematical operations
\usepackage{mathtools}
\usepackage{mathrsfs} % Provides \mathscr command
\usepackage{xfrac} % renders diagonal frac notation: use \sfrac{}{}
\usepackage{dsfont} % renders '1' for characteristic function...
\usepackage{ifthen}
\usepackage{emptypage}
\usepackage{parskip} % to make vertical spaces between two paragraphs
\setlength{\parindent}{15pt} % adds indentation at the starting of each chapter

\usepackage{epigraph} % Epigraphs are the pithy quotations often found at the start (or end) of a chapter
\usepackage{lipsum}

\usepackage[
    a4paper,
    bindingoffset=0.2in,
    centering,
    marginparwidth=2in,
    textwidth=6.1in,
    marginparsep=2em,
    top=2.5cm,
    bottom=2.5cm,
%    showframe
]{geometry} 

%%
%%%Header styles
%%
\usepackage{fancyhdr}
\pagestyle{fancy}
\fancyhf{}
\fancyhead[RE]{\textbf{\textit{\nouppercase{\leftmark}}}}
\fancyhead[LO]{\textbf{\textit{\nouppercase{\rightmark}}}}
\fancyfoot[C]{\thepage} 
\renewcommand{\headrulewidth}{2pt}
\renewcommand{\footrulewidth}{0pt}

%%
%%% Fancy Chapter Styles
%%
%\usepackage[Bjornstrup]{fncychap}
%Options: Conny, Bjarne, Glenn, Lenny, Sonny, Rejne, Bjornstrup

% Custom chapter heading
\makeatletter
\def\thickhrulefill{\leavevmode \leaders \hrule height 1ex \hfill \kern \z@}
\def\@makechapterhead#1{%
  \vspace*{10\p@}%
  {\parindent \z@ \raggedleft \reset@font
    \scshape \@chapapp{} \thechapter
    \par\nobreak
    \interlinepenalty\@M
    \Huge \bfseries #1\par\nobreak
    \hrulefill
    \par\nobreak
    \vskip 40\p@ % <- Reduced from 100pt to 40pt
  }}
\def\@makeschapterhead#1{%
  \vspace*{10\p@}%
  {\parindent \z@ \raggedleft \reset@font
    \scshape \vphantom{\@chapapp{} \thechapter}
    \par\nobreak
    \interlinepenalty\@M
    \Huge \bfseries #1\par\nobreak
    \hrulefill
    \par\nobreak
    \vskip 40\p@ % <- Reduced from 100pt to 40pt
  }}
\makeatother



%%%%%%%%%%%%%%%%%%%%%%%%%%%%%%%%%%%%%%%%%%%%%%%%%%%%%%%%%%%%%%%%%%%%%%%
%%%%%%%%%%%%%%%%%%%%%%%%%%%%%%%%%%%%

%%----------------------------------------------------
%%		 Theorem Styles
%%-----------------------------------------------------
%%%%
%%%% The next few commands set up the theorem type environments.
%%%% Here they are set up to be numbered section.number, but this can
%%%% be changed.
%%%%

\theoremstyle{plain}% default

\newtheorem{theorem}{Theorem}[chapter]
\newtheorem{lem}[theorem]{Lemma}
\newtheorem{prop}{Proposition}[chapter]
\newtheorem{cor}[theorem]{Corollary} % '*' produces an output without numbering

%%
%% If some other type is need, say conjectures, then it is constructed
%% by editing and uncommenting the following.
%%

%\newtheorem{conj}[theorem]{Conjecture} 


%%% 
%%% The following gives definition type environments (which only differ
%%% from theorem type environments in the choices of fonts).  The
%%% numbering is still tied to the theorem counter.
%%% 

\theoremstyle{definition}

\newtheorem{definition}{Definition}[chapter]
\newtheorem{example}{Example}[chapter]
\newtheorem{exercise}[example]{Exercise}

%%
%% Again more of these can be added by uncommenting and editing the
%% following. 
%%

\newtheorem{axiom}[definition]{Axiom}
%\newtheorem{condition}{Condition}
%\newtheorem{hypothesis}[definition]{Hypothesis}
\newtheorem{question}[definition]{Question}


%%% 
%%% The following gives remark type environments (which only differ
%%% from theorem type environments in the choices of fonts).  The
%%% numbering is still tied to the theorem counter.
%%% 

\theoremstyle{remark}

\newtheorem{remark}{Remark}[chapter]
\newtheorem*{note}{Note}
\newtheorem{case}{Case}

%% Add more here 
\newtheorem*{notation}{Notation}
%\newtheorem*{summ}{Summary}
%\newtheorem*{acknow}{Acknowledgment}

\renewcommand\qedsymbol{$\blacksquare$} % redefining the qed symbol

%%%
%%% The following, if uncommented, numbers equations within sections.
%%% 
\numberwithin{equation}{section}

%%
%%% Setting Tcolorbox for boxed 'theorems'
%%
\usepackage{tikz-cd} % Online editor: https://tikzcd.yichuanshen.de/

\usepackage[most]{tcolorbox}
% Usages: \begin{tcolorbox} Here is some text! \end{tcolorbox}

\newtcbtheorem{theo}{Theorem}{}{theorem}
%  \begin{theo}{Name} My theorem  statement \end{theo}
 

%%%%%%%%%% Bibliography management %%%%%%%%%%
\usepackage[titletoc]{appendix}

\usepackage[english]{babel}
\usepackage{csquotes}
\usepackage[nottoc]{tocbibind}
\usepackage[backend=bibtex,style=alphabetic,sorting=ynt]{biblatex}
\addbibresource{bibliography.bib}


%%% Hyper reference set up
\usepackage[bookmarksopen,bookmarksnumbered]{hyperref}
\hypersetup{
	pdftitle={\pdfTitle},
	pdfauthor={\pdfAuthor},
	pdfsubject={\pdfSubject},
	pdfcreationdate={\pdfCreationDate},
	pdfcreator={\pdfCreator},
	pdfkeywords={\pdfKeywords},
	colorlinks=\pdfColorLink,
	linkcolor={\pdfLinkColor},
	%    filecolor=magenta,      
	urlcolor=\pdfUrlColor,
	citecolor=\pdfCiteColor,
	pdfpagemode=UseOutlines,
}
\usepackage{bookmark}

\usepackage{imakeidx}
\makeindex[program=makeindex,options=-s style.ist,columns=2,intoc=true]

%%%%% Custom Commands %%%%%%%%%%%%%%
%% -----------------------------------------------------
%% Custom Macros for LaTeX
%% -----------------------------------------------------
%% Author: Indrajit Ghosh
%% Homepage: https://www.indrajitghosh.onrender.com
%% This file contains user-defined macros to enhance LaTeX functionality.
%% -----------------------------------------------------

% ########## Without math mode ##########

% Defines the Matlab symbol in small caps
\newcommand{\matlab}{{\sc Matlab} }

% Defines Holder's inequality with a properly formatted "o" character
\newcommand{\HI}{H$\ddot{\text{o}}$lder's inequality }

% Macro for bold text: usage \tb{text}
\newcommand{\tb}[1]{\textbf{#1}}

% Macro for normal text: usage \tn{text}
\newcommand{\tn}[1]{\textnormal{#1}}


% ########## Inside Math Mode ##########

% -------- Set Theory --------
% Right arrow symbol
\newcommand{\ra}{\rightarrow}

% -------- Calculus Symbols --------
% Bold vector notation
\newcommand{\cvec}[1]{{\mathbf #1}}

% Bold vector with arrow notation
\newcommand{\rvec}[1]{\vec{\mathbf #1}}

% Unit vector notation
\newcommand{\ihat}{\hat{\textbf{\i}}}
\newcommand{\jhat}{\hat{\textbf{\j}}}
\newcommand{\khat}{\hat{\textbf{k}}}

% Divergence operator
\newcommand{\mdiv}{{\rm div}}

% Gradient and curl operators
\newcommand{\grad}{\textnormal{grad}~}
\newcommand{\curl}{\textnormal{curl}~}

% -------- Linear Algebra --------
% Minor and trace functions
\newcommand{\minor}{{\rm minor}}
\newcommand{\trace}{{\rm trace}}

% Span function
\newcommand{\spn}{{\rm Span}}

% Null space, domain, and range
\newcommand{\nul}{\textnormal{\textrm{null}}}
\newcommand{\dom}{\textnormal{\textrm{dom}}}
\newcommand{\ran}{{\rm ran}}

% Homomorphism
\newcommand{\Hom}{{\rm{Hom}}}

% m x n matrices
\newcommand{\mbyn}[2]{#1\times #2}

% Tensor product and isomorphism
\newcommand{\tensor}{{\otimes}}
\newcommand{\iso}{{\cong}}

% -------- Analysis --------
% Common number sets
\newcommand{\C}{\mathbb{C}}
\newcommand{\R}{\mathbb{R}}
\newcommand{\N}{\mathbb{N}}
\newcommand{\Q}{\mathbb{Q}}
\newcommand{\Z}{\mathbb{Z}}
\newcommand{\Sp}{\mathbb{S}}
\newcommand{\F}{\mathbb{F}} % Field F

% Inner product and empty set
\newcommand{\<}{\langle}
\renewcommand{\>}{\rangle}
\renewcommand{\emptyset}{\varnothing}

% Specific vector spaces
\newcommand{\Cn}{\C^n}
\newcommand{\Zn}{\Z^n}

% -------- Operator Algebra --------
% Bounded operators and matrix algebra
\newcommand{\Bh}{\mathcal{B}(\mathcal{H})}
\newcommand{\Ch}{\mathcal{C}(\mathcal{H})}
\newcommand{\mnc}{\mathbb{M}_n(\mathbb{C})}
\newcommand{\mn}[1]{\mathbb{M}_{#1}(\mathbb{C})}

% von Neumann algebra and convergence
\newcommand{\RR}{\mathscr{R}}
\newcommand{\wkly}{\rightharpoonup}

% Norm symbols
\newcommand{\norm}[1]{\left| \left| #1 \right| \right|}
\newcommand{\opnorm}[1]{\left| \left| #1 \right| \right|_{\text{op}}}
\newcommand{\supnorm}[1]{\left| \left| #1 \right| \right|_{\text{$\infty$}}}

% Direct sum
\newcommand{\dsum}[2]{\bigoplus_{#1}^{#2}}

% Projection relationships
\newcommand{\weaker}{\precsim}
\newcommand{\strictlyweaker}{\prec}

% Murray-von Neumann equivalence
\newcommand{\mvnequiv}[1]{\stackrel{\mathclap{\normalfont\tiny\mbox{$#1$}}}{\sim}}

% -------- Mathematical Operators --------
% Distance function
\DeclareMathOperator{\dist}{dist}

% Floor and ceiling functions
\newcommand{\floor}[1]{\left\lfloor #1 \right\rfloor}
\newcommand{\ceil}[1]{\left\lceil #1 \right\rceil}

% Quotations
\newcommand{\quotes}[1]{\textquotedblleft #1\textquotedblright}

% Differential symbols
\newcommand{\diff}{\mathop{}\!d}
\newcommand{\dd}{\mathop{}\!\mathrm{d}}

% Partial derivatives
\newcommand{\pd}[2]{\frac{\partial #1 }{\partial #2}}
\newcommand{\pdS}[2]{\frac{\partial^2 #1 }{\partial #2^2}}
\newcommand{\pdT}[2]{\frac{\partial^2 #1 }{\partial #2^2}}

% Greek letters
\newcommand{\ep}{\varepsilon}

% Fraktur, calligraphy, and math script
\newcommand{\fA}{\mathfrak A}     \newcommand{\sA}{\mathcal A}		\newcommand{\kA}{\mathscr{A}}
\newcommand{\fB}{\mathfrak B}     \newcommand{\sB}{\mathcal B}		\newcommand{\kB}{\mathscr{B}}
\newcommand{\fC}{\mathfrak C}     \newcommand{\sC}{\mathcal C}		\newcommand{\kC}{\mathscr{C}}
\newcommand{\fD}{\mathfrak D}     \newcommand{\sD}{\mathcal D}		\newcommand{\kD}{\mathscr{D}}
\newcommand{\fE}{\mathfrak E}     \newcommand{\sE}{\mathcal E}		\newcommand{\kE}{\mathscr{E}}
\newcommand{\fF}{\mathfrak F}     \newcommand{\sF}{\mathcal F}		\newcommand{\kF}{\mathscr{F}}
\newcommand{\fG}{\mathfrak G}     \newcommand{\sG}{\mathcal G}		\newcommand{\kG}{\mathscr{G}}
\newcommand{\fH}{\mathfrak H}     \newcommand{\sH}{\mathcal H}		\newcommand{\kH}{\mathscr{H}}
\newcommand{\fI}{\mathfrak I}     \newcommand{\sI}{\mathcal I}		\newcommand{\kI}{\mathscr{I}}
\newcommand{\fK}{\mathfrak K}     \newcommand{\sK}{\mathcal K}		\newcommand{\kK}{\mathscr{K}}
\newcommand{\fJ}{\mathfrak J}     \newcommand{\sJ}{\mathcal J}		\newcommand{\kJ}{\mathscr{J}}
\newcommand{\fL}{\mathfrak L}     \newcommand{\sL}{\mathcal L}		\newcommand{\kL}{\mathscr{L}}
\newcommand{\fM}{\mathfrak M}     \newcommand{\sM}{\mathcal M}		\newcommand{\kM}{\mathscr{M}}
\newcommand{\fN}{\mathfrak N}     \newcommand{\sN}{\mathcal N}		\newcommand{\kN}{\mathscr{N}}
\newcommand{\fO}{\mathfrak O}     \newcommand{\sO}{\mathcal O}		\newcommand{\kO}{\mathscr{O}}
\newcommand{\fP}{\mathfrak P}     \newcommand{\sP}{\mathcal P}		\newcommand{\kP}{\mathscr{P}}
\newcommand{\fQ}{\mathfrak Q}     \newcommand{\sQ}{\mathcal Q}		\newcommand{\kQ}{\mathscr{Q}}
\newcommand{\fR}{\mathfrak R}     \newcommand{\sR}{\mathcal R}		\newcommand{\kR}{\mathscr{R}}
\newcommand{\fS}{\mathfrak S}     \newcommand{\sS}{\mathcal S}		\newcommand{\kS}{\mathscr{S}}
\newcommand{\fT}{\mathfrak T}     \newcommand{\sT}{\mathcal T}		\newcommand{\kT}{\mathscr{T}}
\newcommand{\fU}{\mathfrak U}     \newcommand{\sU}{\mathcal U}		\newcommand{\kU}{\mathscr{U}}
\newcommand{\fV}{\mathfrak V}     \newcommand{\sV}{\mathcal V}		\newcommand{\kW}{\mathscr{W}}
\newcommand{\fW}{\mathfrak W}     \newcommand{\sW}{\mathcal W}		\newcommand{\kV}{\mathscr{V}}
\newcommand{\fX}{\mathfrak X}     \newcommand{\sX}{\mathcal X}		\newcommand{\kX}{\mathscr{X}}
\newcommand{\fY}{\mathfrak Y}     \newcommand{\sY}{\mathcal Y}		\newcommand{\kY}{\mathscr{Y}}
\newcommand{\fZ}{\mathfrak Z}     \newcommand{\sZ}{\mathcal Z}		\newcommand{\kZ}{\mathscr{Z}}


%bold vector macros
\newcommand{\bu}{{\bf u}}
\newcommand{\bv}{{\bf v}}
\newcommand{\bw}{{\bf w}}
\newcommand{\bx}{{\bf x}}
\newcommand{\by}{{\bf y}}
\newcommand{\bz}{{\bf z}}
\newcommand{\bX}{{\bf X}}
\newcommand{\bY}{{\bf Y}}
\newcommand{\bZ}{{\bf Z}}
\newcommand{\btheta}{{\bf \theta}}
\newcommand{\bTheta}{{\bf \Theta}}
\newcommand{\bze}{{\bf 0}}
\newcommand{\bT}{{\bf T}}
\newcommand{\bB}{{\bf B}}
\newcommand{\bE}{{\bf E}}
\newcommand{\bN}{{\bf N}}


%%%%%%%%%%%%%%%%%%%%%%%%%%%%%%%%%%%%%%%%%%%%%%%%%%%%%%%%%%%%%%%%%%%%%%%%%%
%%%         Document Specific Commands
%%%%%%%%%%%%%%%%%%%%%%%%%%%%%%%%%%%%%%%%%%%%%%%%%%%%%%%%%%%%%%%%%%%%%%%%%%
\newcommand{\irchi}[2]{\raisebox{\depth}{$#1\chi$}} % inner command, used by \rchi
\DeclareRobustCommand{\rchi}{{\mathpalette\irchi\relax}}

\newcommand{\ipdt}[2]{\left\langle {#1, #2} \right\rangle}
\newcommand{\raR}{\textbf{R}}
\newcommand{\nulN}{\textbf{N}}
\newcommand{\graph}[1]{\text{Gr}\left( #1\right)}
\newcommand{\krein}[1]{{#1}^{\textbf{Kr}}}
\newcommand{\fried}[1]{{#1}^{\textbf{Fr}}}
\newcommand{\proj}[1]{#1_\tn{proj}} % projection lattice of an algebra
\newcommand{\sa}[1]{#1_{\tn{sa}}} % self-adjoint part of an algebra

\newcommand{\afftext}{\mathrm{aff}}
\newcommand{\aff}[1]{#1_{\afftext}}
\newcommand{\mvnaff}[1]{#1_{\afftext}^{\mathrm{MvN}}}
\newcommand{\affc}[1]{#1_{\afftext}^{c}}
\newcommand{\affs}[1]{\mathrm{Aff}_{s}\left( #1 \right)}
\newcommand{\waff}[1]{#1_{\afftext}^{\text{wor}}}
\newcommand{\Lim}{{\textnormal{\textsc{Lim}}} ~}
\newcommand{\spec}[1]{\textnormal{sp}\left( #1 \right)}

\newcommand{\cstarcat}{\mathbf{C^*Alg}}
\newcommand{\mstarcat}{\mathbf{M\sigma C^*Alg}}
\newcommand{\wstarcat}{\mathbf{W^*Alg}}

\newcommand{\mat}[2]{\mathbb{M}_#1 \left( #2 \right)}

\newcommand{\invran}{\mathrel{\overset{\makebox[0pt]{\mbox{\normalfont\tiny\sffamily inv}}}{\circ}}}

\newcommand{\blah}{{\textcolor{red}{~blah ??}}}
\newcommand{\inprogress}{{\textcolor{magenta}{\centering \LARGE \texttt{... Chapter in progress ...} \\}}}



%%%%%%%%%%% Glossaries and Acronyms %%%%%%%%%%%%%
\usepackage[%
acronym,%
symbols,%
nogroupskip,%
nonumberlist,% Uncomment for not printing the page numbers at glossaries
%sort=use
]%
{glossaries-extra}
\makenoidxglossaries
% !TEX root = ../main.tex
% =========================================================================
% symbols.tex
% Template for defining mathematical symbols for the glossary
% Author: Indrajit Ghosh
% Created On: Oct 21, 2025
% =========================================================================
%
% Usage:
%  - Define a symbol: \glsxtrnewsymbol[sort={number},description={Description}]{label}{LaTeX code}
%  - Use the symbol in text or math via: \gls{label}
%    Example: $\gls{R}$ or \gls{Bh}
%
% Notes:
%  - sort={number} determines the order in the List of Symbols
%  - description={...} is what appears in the List of Symbols
%  - label is the key used with \gls{}
%  - LaTeX code is what will be rendered in the document

% ---------------------------
% Basic Number Sets
% ---------------------------
\glsxtrnewsymbol[sort={1},description={Set of all natural numbers}]{N}{\ensuremath{\N}}
\glsxtrnewsymbol[sort={2},description={Set of all integers}]{Z}{\ensuremath{\Z}}
\glsxtrnewsymbol[sort={3},description={Set of all real numbers}]{R}{\ensuremath{\R}}
\glsxtrnewsymbol[sort={4},description={Field of complex numbers}]{C}{\ensuremath{\C}}

% ---------------------------
% Matrices and Operators
% ---------------------------
\glsxtrnewsymbol[sort={5},description={$n\times n$ matrix algebra over $\mathbb{C}$}]{mnc}{\ensuremath{\mnc}}
\glsxtrnewsymbol[sort={6},description={Algebra of bounded operators on a Hilbert space $\mathcal{H}$}]{Bh}{\ensuremath{\Bh}}
\glsxtrnewsymbol[sort={7},description={Closed unbounded operators on $\mathcal{H}$}]{Ch}{\ensuremath{\Ch}}
\glsxtrnewsymbol[sort={8},description={Graph of an operator $T$}]{grapht}{\ensuremath{\graph{T}}}

% ---------------------------
% Special Operators and Projections
% ---------------------------
\glsxtrnewsymbol[sort={9},description={Kaufman inverse of operator $T$}]{tdag}{\ensuremath{T^\dagger}}
\glsxtrnewsymbol[sort={10},description={Characteristic projection of the closed operator $T$}]{chi}{\ensuremath{\rchi(T)}}
\newcommand{\1}[1]{\mathds{1}_{#1}} % characteristic function
\glsxtrnewsymbol[sort={11},description={Indicator function of the set $\Omega$}]{indicator}{\ensuremath{\1{\Omega}}}

% ---------------------------
% von Neumann Algebra Symbols
% ---------------------------
\glsxtrnewsymbol[sort={12},description={Self-adjoint part of the von Neumann algebra $\mathscr{M}$}]{msa}{\ensuremath{\sa{\kM}}}
\glsxtrnewsymbol[sort={13},description={Lattice of projections in $\mathscr{M}$}]{projm}{\ensuremath{\proj{\kM}}}
\glsxtrnewsymbol[sort={14},description={Lattice of affiliated subspaces of $\mathscr{M}$}]{affsm}{\ensuremath{\affs{\kM}}}
\glsxtrnewsymbol[sort={15},description={Murray–von Neumann affiliated operators associated with $\mathscr{M}$}]{mvnaff}{\ensuremath{\mvnaff{\kM}}}
\glsxtrnewsymbol[sort={16},description={Densely-defined affiliated operators}]{ddaff}{\ensuremath{\aff{\kM}^{dd-c}}}
\glsxtrnewsymbol[sort={17},description={All affiliated operators associated with $\mathscr{M}$}]{maff}{\ensuremath{\aff{\kM}}}
\glsxtrnewsymbol[sort={18},description={Affiliated operators that are closed}]{maffc}{\ensuremath{\affc{\kM}}}

% ---------------------------
% C*-Algebra Symbols
% ---------------------------
\glsxtrnewsymbol[sort={19},description={Affiliated operators associated with a $C^*$-algebra $\fA$}]{aaff}{\ensuremath{\aff{\fA}}}
\glsxtrnewsymbol[sort={20},description={$n\times n$ matrices over a $C^*$-algebra $\fA$}]{mna}{\ensuremath{\mat{n}{\fA}}}

% ---------------------------
% Categories and Extensions
% ---------------------------
\glsxtrnewsymbol[sort={21},description={Category of unital $C^*$-algebras with unital *-homomorphisms}]{cstar}{\ensuremath{\cstarcat}}
\glsxtrnewsymbol[sort={22},description={Category of von Neumann algebras with normal *-homomorphisms}]{wstar}{\ensuremath{\wstarcat}}
\glsxtrnewsymbol[sort={23},description={Krein extension of the operator $S$}]{skr}{\ensuremath{\krein{S}}}
\glsxtrnewsymbol[sort={24},description={Friedrich extension of the operator $S$}]{sfr}{\ensuremath{\fried{S}}}

% ---------------------------
% Add new symbols here!
% Example:
% \glsxtrnewsymbol[sort={25},description={Your description here}]{newlabel}{\ensuremath{symbol_code}}

% =========================================================================
% Optional: include all symbols in the List of Symbols
% Remove the following line if you want only referenced symbols to appear
\glsaddall


\usepackage{acro}
% !TEX root = ../main.tex
% Abbreviations
% You can use these abbreviations in the thesis using the cmd '\ac{}'
% For example, \ac{masa}

\DeclareAcronym{masa}{
  short=MASA,
  long=Maximal Abelian Self-adjoint Algebras,
}
\DeclareAcronym{vna}{
  short=vNa,
  long=von Neumann Algebra,
}
\DeclareAcronym{mvn}{
  short=MvN,
  long=Murray-von Neumann,
}



%%%%%%%%%%%%%%%%%%%%%%%%%%%%%%%%%%%%%%%%%%%%%%%%%%%%%%




%%
%% This is the end of the preamble.
%%

    


% =========================================================================
% USAGE GUIDE: Theorem-Like Environments
% =========================================================================
%
% Your preamble defines standard AMS-style environments and some additional
% helpers for structuring mathematical content. Use them as follows:
%
% -------------------
% Theorem Environments
% -------------------
%
% \begin{theorem}[Optional Name]
%   Statement of the theorem goes here.
% \end{theorem}
%
% \begin{lem}[Optional Name]
%   Statement of the lemma.
% \end{lem}
%
% \begin{prop}[Optional Name]
%   Statement of the proposition.
% \end{prop}
%
% \begin{cor}[Optional Name]
%   Statement of the corollary.
% \end{cor}
%
% -------------------
% Definition/Example
% -------------------
%
% \begin{definition}[Optional Name]
%   Define a concept, object, or structure.
% \end{definition}
%
% \begin{example}[Optional Name]
%   Provide an example illustrating a definition or theorem.
% \end{example}
%
% \begin{exercise}[Optional Name]
%   Pose a problem or exercise for the reader (optional numbering with example counter).
% \end{exercise}
%
% -------------------
% Remarks/Notes
% -------------------
%
% \begin{remark}[Optional Name]
%   Add remarks or observations.
% \end{remark}
%
% \begin{note}
%   Use for unnumbered notes.
% \end{note}
%
% \begin{case}[Optional Name]
%   Use for enumerating cases in proofs.
% \end{case}
%
% \begin{notation}
%   State notational conventions.
% \end{notation}
%
% -------------------
% tcolorbox (Optional Boxed Theorems)
% -------------------
%
% For visually highlighted statements, use tcolorbox with theorem styling:
%
% \begin{theo}{Name of Theorem}{label}
%   Statement of theorem goes here.
% \end{theo}
%
% This produces a boxed theorem with optional label for referencing.
%
% =========================================================================
% NOTES:
% - All theorem-like environments are numbered by chapter by default.
% - Equations are numbered within sections (see preamble: \numberwithin{equation}{section}).
% - Optional names for theorems/lemmas appear in parentheses after the environment.
% =========================================================================


\begin{document}

    % Coverpage
\begin{titlepage}
    \begin{center}

\vspace*{8ex}
{\huge \bfseries \ThesisTitle}\\[0.4cm]
\ifthenelse{\equal{\ThesisSubtitle}{}}{}{%
    {\Large \textit{\ThesisSubtitle}}\\[5ex]
}

\vspace{14ex}
\href{\AuthorHomepage}{\Large\scshape \Author} %

\vspace{18ex}
\includegraphics[width=40mm]{images/isi_logo.png} \\[1cm] 

{\LARGE \Institute}\\[1.5cm]

{\large \ThesisYear}\\[1cm] % Date


\end{center}
\end{titlepage}

    \thispagestyle{empty}
    \clearpage

    \frontmatter

    % Title page for the thesis.
\begin{titlepage}
	\begin{center}
		
		\vspace{.06\textheight}
		{\scshape\LARGE \Institute \par}\vspace{1.5cm} 
		\textsc{\Large Doctoral Thesis}\\[0.5cm] 

		\hrulefill \\[0.4cm]

		{\huge \bfseries \textcolor{indraBlue}{\ThesisTitle}\par}\vspace{0.2cm}
        \ifthenelse{\equal{\ThesisSubtitle}{}}{}{%
            {\Large \textit{\textcolor{indraBlue}{\ThesisSubtitle}}\par}\vspace{0.4cm}
        }
		\hrulefill \\[0.4cm]

		\begin{minipage}[t]{0.4\textwidth}
			\begin{flushleft} \large
				\emph{Author:}\\
				{\textcolor{indraBlue}{\Author}} 
			\end{flushleft}
		\end{minipage}
		\begin{minipage}[t]{0.4\textwidth}
			\begin{flushright} \large
				\emph{Advisor:} \\
				{\textcolor{indraBlue}{\PhdAdvisor}} 
			\end{flushright}
		\end{minipage}\\[1cm]

		\vfill

		\large \textit{A thesis submitted to the Indian Statistical Institute \\in partial fulfilment of the requirements for\\ the degree of \\
        Doctor of Philosophy (in Mathematics)}\\[1cm]

		\vfill

		\Large \Department\\[.25cm] \Institute
		\ifthenelse{\equal{\Centre}{}}{}{, \Centre}
		\\[2cm]

		{\large \ThesisYear}\\[4cm] 

		\vfill
	\end{center}
\end{titlepage}

    \thispagestyle{empty}

    \chapter*{Abstract}
    % chapters/abstract.tex
% !TEX root = ../main.tex
%
% Template Abstract for PhD Thesis LaTeX Template
% Author: Indrajit Ghosh
% Created: 2025-10-21

This LaTeX file was originally written while I was completing my PhD at the Indian Statistical Institute, Bangalore. After defending my thesis, I decided to convert it into a general-purpose template that can be used by anyone writing a thesis in LaTeX, and especially by ISI PhD scholars. The template includes a structured preamble, theorem environments, macros, color schemes, symbol and acronym management, and sample chapters.  

It is designed to be flexible and easily adaptable, allowing users to focus on writing their thesis content without worrying about formatting. By following the examples in this template, scholars can create professional-looking theses with consistent styling, automated lists of symbols and acronyms, and well-formatted mathematical content.  

This template aims to provide a practical starting point for LaTeX-based thesis writing while showcasing best practices for organizing large LaTeX projects.

    \addcontentsline{toc}{chapter}{Abstract}

    \chapter*{Dedication}
    % !TEX root = ../main.tex

This template is dedicated to \LaTeX,  
for faithfully not crashing (most of the time),  
for making my equations look nicer than my handwriting,  
and for being the silent companion of many sleepless thesis nights.  

Also dedicated to anyone who has ever spent hours debugging a missing brace or forgotten \verb|\end{document}|.  
May your compilations be fast and your error logs short!
    \addcontentsline{toc}{chapter}{Dedication}

    \chapter*{Acknowledgements}
    % !TEX root = ../main.tex

First and foremost, I thank \LaTeX\ for being a steadfast companion through countless hours of writing, compiling, and debugging. Without it, this template would probably still be a pile of loose notes and coffee stains.  

I also extend my gratitude to spellcheckers, autocorrect, and online forums—especially Stack Exchange—for rescuing me from typos, mysterious compilation errors, and the occasional existential crisis caused by a missing brace.  

Finally, thanks to anyone who ever glanced at this template and thought, “Hey, this might actually save me some time!” May your citations be complete, your figures properly aligned, and your bibliography entries forever consistent.

\vspace{1cm}
\begin{flushright}
\textit{--- Indrajit Ghosh}
\end{flushright}

    \addcontentsline{toc}{chapter}{Acknowledgements}

    \pdfbookmark[section]{\contentsname}{toc}
    \tableofcontents

    \printnoidxglossary[type=symbols,style=long,title={List of Symbols}]
    \thispagestyle{empty}

    \mainmatter

    % =============================================================
    % Funny Mathematical Chapters
    % =============================================================

    \chapter{Introduction}
    \label{chap:intro}
    
\epigraph{"Mathematics is serious, but your LaTeX doesn't have to be."}{\textit{Anonymous TeX Enthusiast}}

Welcome to this templated PhD thesis! This chapter is your **user manual**, showing how to wield the magical powers of LaTeX, AMS environments, colors, macros, symbols, and acronyms — all in one place.

% ===========================================================
\section{Files in This Template}\index{template}
This template consists of several supporting files, each with a specific role:

\begin{itemize}
    \item \textbf{main.tex:} The master document that ties everything together. Compile this to produce the final PDF.
    \item \textbf{preamble.tex:} Loads packages, defines theorem environments, AMS math structures, headers, and page styles.
    \item \textbf{colors.tex:} Defines all custom colors. Use them via \verb|\textcolor{colorname}{text}| or in hyperref, tikz, and tcolorbox.
    \item \textbf{macros.tex:} Handy macros for math and text. Examples:
        \begin{itemize}
            \item \verb|\R| $\Rightarrow \mathbb{R}$, \verb|\C| $\Rightarrow \mathbb{C}$
            \item \verb|\Bh| $\Rightarrow \mathcal{B}(\mathcal{H})$
            \item \verb|\tb{text}| $\Rightarrow$ \textbf{text}, \verb|\tn{text}| $\Rightarrow$ normal text
            \item Vector macros: \verb|\bv|, \verb|\bx|, \verb|\btheta|, etc.
        \end{itemize}
    \item \textbf{chapters/}: Contains separate `.tex` files for each chapter.
    \item \textbf{chapters/symbols.tex:} Defines all mathematical symbols for the glossary. Use them via \verb|\gls{}|, e.g., $\gls{R}$.
    \item \textbf{chapters/acronyms.tex:} Defines all abbreviations. Use via \verb|\ac{}|, e.g., \ac{masa}.
\end{itemize}

% ===========================================================
\section{How to Use Theorem-Like Environments}
Thanks to \verb|preamble.tex|, you have access to **AMS theorem environments** and boxed theorems.  

\subsection*{Theorems, Lemmas, Propositions, Corollaries}
\begin{itemize}
    \item \verb|\begin{theorem}[Optional Name] ... \end{theorem}|
    \item \verb|\begin{lem}[Optional Name] ... \end{lem}|
    \item \verb|\begin{prop}[Optional Name] ... \end{prop}|
    \item \verb|\begin{cor}[Optional Name] ... \end{cor}|
\end{itemize}

\subsection*{Definitions, Examples, Exercises}
\begin{itemize}
    \item \verb|\begin{definition}[Optional Name] ... \end{definition}|
    \item \verb|\begin{example}[Optional Name] ... \end{example}|
    \item \verb|\begin{exercise}[Optional Name] ... \end{exercise}|
\end{itemize}

\subsection*{Remarks, Notes, Cases, Notation}
\begin{itemize}
    \item \verb|\begin{remark}[Optional Name] ... \end{remark}|
    \item \verb|\begin{note} ... \end{note}| (unnumbered)
    \item \verb|\begin{case}[Optional Name] ... \end{case}|
    \item \verb|\begin{notation} ... \end{notation}|
\end{itemize}

\subsection*{Boxed Theorems}
\verb|tcolorbox| provides visually appealing boxed theorems:
\begin{verbatim}
\begin{theo}{Name of Theorem}{label}
  Statement goes here.
\end{theo}
\end{verbatim}

% ===========================================================
\section{Symbols and Glossary}\index{symbols}
All symbols are defined in \verb|chapters/symbols.tex|.  

Use the command \verb|\gls{}| to reference them. Examples:

\begin{itemize}
    \item $\gls{N}$: Set of natural numbers  
    \item $\gls{R}$: Set of real numbers  
    \item $\gls{Bh}$: Algebra of bounded operators on a Hilbert space  
    \item $\gls{affsm}$: Lattice of affiliated subspaces of a von Neumann algebra  
    \item $\gls{mvnaff}$: Murray–von Neumann affiliated operators  
\end{itemize}

\begin{tcolorbox}[colback=goldenrod!10,colframe=crimson,title=Tip]
Always use \verb|\gls{}| instead of typing symbols manually — this ensures they appear in the \textbf{List of Symbols}.
\end{tcolorbox}

% ===========================================================
\section{Acronyms}
Defined in \verb|chapters/acronyms.tex|.  

Use \verb|\ac{}| to reference them. Examples:

\begin{itemize}
    \item \ac{masa}: Maximal Abelian Self-adjoint Algebra  
    \item \ac{vna}: von Neumann Algebra  
    \item \ac{mvn}: Murray–von Neumann  
\end{itemize}

\begin{tcolorbox}[colback=softblue!10,colframe=royalblue,title=Tip]
Introduce acronyms at first use using \verb|\ac{}|, then simply repeat using \verb|\ac{}|. Your readers (and future self) will thank you.
\end{tcolorbox}

% ===========================================================
\section{AMS and Math Environments}
All AMS environments from \verb|amsmath, amssymb, mathtools| are available:
\begin{itemize}
    \item Inline math: \verb|$...$|
    \item Display math: \verb|\[ ... \]|
    \item Alignments: \verb|align, gather, multline, cases|
    \item Numbering: Equations are numbered per section (\verb|\numberwithin{equation}{section}|)
\end{itemize}

% ===========================================================
\section{Colors and Macros}
\begin{itemize}
    \item Colors: \verb|colors.tex| provides \textcolor{royalblue}{royalblue}, \textcolor{indraRed}{indraRed}, \textcolor{deepblue}{deepblue}, etc.
    \item Macros: \verb|macros.tex| provides shortcuts for math and text. Examples:
    \begin{itemize}
        \item Vectors: \verb|\bv|, \verb|\bx|, \verb|\btheta|
        \item Operators: \verb|\grad|, \verb|\curl|, \verb|\div|
        \item Algebras: \verb|\Bh|, \verb|\Ch|, \verb|\aff{A}|
        \item Bold or normal text: \verb|\tb{text}|, \verb|\tn{text}|
    \end{itemize}
\end{itemize}

% ===========================================================
\section{Compiling the Thesis}
\begin{enumerate}
    \item Always compile \verb|main.tex|.
    \item Make sure \verb|preamble.tex|, \verb|colors.tex|, and \verb|macros.tex| are accessible via \verb|\input{}|.
    \item To generate the glossary and symbols:
    \begin{itemize}
        \item \verb|makeglossaries main|
        \item \verb|makeindex main|
    \end{itemize}
    \item View your PDF to check formatting, symbols, and theorem boxes.
\end{enumerate}

% ===========================================================
\section{Final Notes}
This chapter is your **LaTeX cheat sheet** for this thesis template. Modify colors, macros, or environments freely — even serious mathematicians deserve a little fun.


    \chapter{Category Theory and the Existential Crisis of Sets}
    \label{chap:category_crisis}
    \epigraph{"Every set is lonely until someone defines a function from it."}{\textit{Anonymous Category Theorist}}

\section{Sets in Therapy}
A set is just a collection of objects that never talk to each other. 
Category theory, in a surprising twist, gives them morphisms so they can at least make small talk.

\begin{definition}
A \emph{category} $\mathcal{C}$ is a structure consisting of:
\begin{itemize}
    \item Objects (with commitment issues),
    \item Morphisms between them (sometimes called ``maps of destiny''),
    \item Composition laws (that rarely commute with breakfast).\index{breakfast}
\end{itemize}
\end{definition}

\begin{theorem}[Yoneda’s Existential Crisis]
Every object is completely defined by how others perceive it. 
Formally, for each $A$ in $\mathcal{C}$,
\[
\mathrm{Hom}_{\mathcal{C}}(-,A)
\]
contains all the gossip about $A$.
\end{theorem}

\begin{proof}
This proof is left as an exercise in self-reflection.
\end{proof}

\begin{remark}
Set theorists often call category theory “philosophy with diagrams.” 
Category theorists call set theory “philosophy without style.”
\end{remark}


    \chapter{The Spectrum Always Rings Twice: Adventures in Functional Analysis}
    \label{chap:spectrum_fun}
    \epigraph{"Every operator has a spectrum; only some are worth listening to."}{\textit{Spectral Analyst, ca. 2025}}

\section{Prelude to Spectral Analysis}
Given a Hilbert space $\mathcal{H}$ and a bounded operator $T \in B(\mathcal{H})$, 
the \emph{spectrum} $\sigma(T)$ is the set of complex numbers $\lambda$ for which $T - \lambda I$ 
is not invertible. 
In other words, it’s where analysis stops and magic begins.

\begin{tcolorbox}
\textbf{Moral:} The resolvent set is like your to-do list — everything outside it is under control.
\end{tcolorbox}

\begin{theorem}[Spectral Decomposition, Coffee Edition]
Every normal operator $T$ can be written as
\[
T = \int_{\sigma(T)} \lambda \, dE(\lambda)
\]
where $E$ is a projection-valued measure, and the integral is taken with sufficient caffeine.
\end{theorem}

\begin{proof}
The proof involves standard techniques from measure theory and late-night panic.
\end{proof}

\begin{example}
If $T$ is self-adjoint, $\sigma(T) \subset \mathbb{R}$.  
If $T$ is a PhD student, $\sigma(T)$ includes several imaginary parts.
\end{example}


    \backmatter

    \nocite{*}
    \printbibliography[
        heading=bibintoc,
        title={Bibliography}]

    \printindex

\end{document}
